\documentclass[letterpaper]{article}
\usepackage[margin=1in]{geometry}
\usepackage{amsmath,amssymb}
\usepackage{booktabs}
\usepackage{multirow}
\usepackage{graphicx}
\usepackage{xcolor}
\usepackage{hyperref}
\usepackage{natbib}
\usepackage{array}
\usepackage{caption}

\title{The Hidden Second Factor: Model$\times$Item Interaction Accounts for 34\% of LLM Evaluation Variance}

\author{Anonymous}
\date{}

\begin{document}
\maketitle

% ============================================================
% ABSTRACT
% ============================================================
\begin{abstract}
When LLM evaluation results vary across reproductions, practitioners cannot determine whether the variance originates from prompt phrasing, random seeds, numerical precision, temperature, or item selection---let alone how these factors interact.
We conduct a multi-benchmark Generalizability Theory (G-theory) study crossing infrastructure and methodology factors across 4 open-weight models (7--9B parameters) on MMLU, GSM8K, ARC-Challenge, and HellaSwag, totaling over 2.3 million evaluation records.
On MMLU (345,600 balanced cross-model records), item selection accounts for 42.1\% of variance and model$\times$item interaction emerges as the second-largest source at 33.6\%, revealing that item difficulty rankings are substantially model-specific: single-model analyses inflate item dominance from 42\% to 61\% by conflating universal and model-specific difficulty.
A Decision study (D-study) predicted the 4-model generalizability coefficient within $\Delta=0.004$ of the observed value, providing out-of-sample validation.
Across four benchmarks, variance structure is approximately invariant (item dominance range = 5.7 percentage points), enabling cautious cross-benchmark extrapolation of MMLU-derived evaluation budgets with a 3--6 percentage-point safety margin.
G-theory requires 1.6--4.1$\times$ more items than IRT for reliable measurement, with the ratio being benchmark-dependent---resolving conflicting conclusions in the literature about IRT sufficiency.
A partial two-GPU study reveals that hardware$\times$methodology interaction reaches 4.2\%, breaking the separability assumption that holds within single-GPU environments ($<$0.1\%).
We release an open-source G-theory variance decomposition pipeline and D-study budget calculator.
\end{abstract}

% ============================================================
% 1. INTRODUCTION
% ============================================================
\section{Introduction}

Reproducibility failures in LLM evaluation undermine the credibility of benchmark-based comparisons.
A model declared superior on MMLU under one evaluation protocol may lose its advantage under a different prompt template \citep{sclar2024quantifying}, random seed \citep{bui2025seed}, or numerical precision setting \citep{yuan2025precision}.
Yet each prior study isolates at most one or two variance sources, implicitly assuming that these factors do not interact.
The field lacks a unified framework for simultaneously attributing evaluation variance to all relevant sources and their interactions.

Generalizability Theory (G-theory; \citealt{shavelson1991generalizability, brennan2001generalizability}) provides exactly this framework.
Originally developed in educational measurement, G-theory decomposes observed variance into contributions from every facet (e.g., items, raters, occasions) and their interactions in a single analysis, then projects reliability under hypothetical measurement designs via Decision studies (D-studies).
\citet{mesting2026generalizability} adapted G-theory to LLM evaluation with five methodological facets on a single benchmark, establishing feasibility but leaving infrastructure factors, cross-benchmark comparisons, and model-as-facet analysis unexplored.

We extend this line of work with three contributions:

\textbf{(1) Model$\times$item interaction quantification.}
In a balanced 4-model MMLU design (345,600 records), model$\times$item interaction accounts for 33.6\% [95\% CI: 28.4--37.9\%] of total variance---the second-largest source after item selection (42.1\%).
Single-model analyses inflate item dominance to 61.1\% by conflating universal and model-specific item difficulty.
A D-study from 3 models predicted the 4-model generalizability coefficient ($G$) within $\Delta = 0.004$, providing out-of-sample validation that G-theory's projective machinery works for LLM evaluation.

\textbf{(2) IRT vs.\ G-theory direct comparison.}
G-theory requires 3.6--4.1$\times$ more items than IRT on MMLU to reach the same reliability target, because it accounts for prompt, seed, and ordering dimensions that IRT ignores.
This ratio is benchmark-dependent, ranging from 1.6$\times$ (HellaSwag) to 4.1$\times$ (MMLU), explaining why the field has not converged on whether IRT-based subsampling is sufficient.

\textbf{(3) Approximate benchmark-invariance of variance structure.}
Across MMLU, GSM8K, ARC-Challenge, and HellaSwag, the item dominance proportion varies by only 5.7 percentage points (38.4--44.1\%), and the variance component ranking is preserved.
MMLU-derived evaluation budgets transfer to other benchmarks with a 3--6 percentage-point $G$-coefficient error, enabling cautious cross-benchmark extrapolation.

% ============================================================
% 2. RELATED WORK
% ============================================================
\section{Related Work}

\paragraph{Variance decomposition in ML evaluation.}
\citet{bouthillier2021accounting} decomposed ML benchmark variance across hyperparameters, data splits, and random seeds via ANOVA, establishing factorial variance decomposition as a methodology for ML evaluation.
\citet{henderson2018deep} documented large variance across seeds in deep reinforcement learning.
\citet{dodge2020finetuning} and \citet{sellam2022multberts} quantified fine-tuning instability from seeds and data ordering in BERT-scale models.
\citet{bui2025seed} extended this to instruction-tuned LLMs, finding extreme seed sensitivity (RTES variance up to 18.22 for RoBERTa on COPA).

\paragraph{Prompt and evaluation sensitivity.}
\citet{sclar2024quantifying} demonstrated up to 76\% accuracy swings from prompt formatting alone.
\citet{lu2022fantastically} and \citet{pezeshkpour2024large} documented option-ordering effects in multiple-choice evaluation.
BrittleBench \citep{brittlebench2026} quantified prompt sensitivity at 50\% variance on fixed hardware.

\paragraph{Infrastructure effects.}
\citet{yuan2025precision} documented up to 9\% accuracy variation from hardware and numerical precision without factorial attribution.
\citet{messina2026floating} identified floating-point non-associativity as a mechanism for cross-hardware divergence.

\paragraph{G-theory in LLM evaluation.}
\citet{mesting2026generalizability} adapted G-theory to LLM evaluation with 5 methodological facets on a single benchmark, finding judge identity as the dominant variance source (43.8\%).
We build on this by (a) adding infrastructure as a facet, (b) treating model as a crossed facet to expose model$\times$item interaction, (c) extending to four benchmarks, and (d) fixing the judge to isolate non-judge variance sources.
\citet{shavelson1991generalizability} and \citet{brennan2001generalizability} provide the canonical G-theory references; \citet{jiang2017generalized} and \citet{mushquash2006generalizability} address G-theory for binary response data.

\paragraph{IRT-based evaluation.}
\citet{lalor2019learning} introduced IRT-based item analysis for NLP evaluation.
\citet{rodriguez2021evaluation} and \citet{vania2021comparing} further developed IRT applications in NLP benchmarking.
\citet{polo2024tinybenchmarks} proposed IRT-based subsampling (tinyBenchmarks), recommending $\sim$100 items for reliable evaluation.
\citet{perlitz2023efficient} explored efficient benchmarking strategies.
Our work provides the first direct comparison of IRT and G-theory item requirements, quantifying the cost of multi-facet reliability.

\paragraph{Alternative frameworks.}
\citet{kunievsky2025alternative} proposed decomposing evaluation variance into intent, articulation, and model uncertainty---a complementary perspective to G-theory's facet-based approach.
\citet{dehghani2021benchmark} argued that benchmark selection itself is a variance source (the ``Benchmark Lottery'').
\citet{biderman2024lessons} synthesized practical lessons for evaluation reproducibility.

% ============================================================
% 3. METHOD
% ============================================================
\section{Method}

\subsection{Generalizability Theory Background}

G-theory \citep{shavelson1991generalizability, brennan2001generalizability} partitions observed score variance into components attributable to each facet and their interactions.
For a fully crossed design with items ($i$), prompts ($p$), seeds ($s$), temperatures ($t$), orderings ($o$), and precision ($r$), the observed score $X_{ipstor}$ decomposes as:
\begin{equation}
X_{ipstor} = \mu + \alpha_i + \beta_p + \gamma_s + \delta_t + \epsilon_o + \zeta_r + (\alpha\beta)_{ip} + \cdots + e_{ipstor}
\end{equation}
where each Greek letter denotes a main effect and parenthesized terms denote interactions.
Variance components $\hat{\sigma}^2$ are estimated via Henderson's Method I (expected mean squares), with REML as a robustness check.

The \textbf{generalizability coefficient} $G$ quantifies reliability:
\begin{equation}
G = \frac{\hat{\sigma}^2_i}{\hat{\sigma}^2_i + \sum_{f \neq i} \frac{\hat{\sigma}^2_f}{n_f}}
\end{equation}
where the sum runs over all non-item variance components divided by their respective sample sizes.

A \textbf{D-study} projects $G$ under hypothetical designs by varying $n_f$, enabling practitioners to determine the number of items, prompts, and seeds needed for a target reliability.

When model is treated as a crossed facet, the item variance component splits into universal item difficulty ($\hat{\sigma}^2_i$) and model$\times$item interaction ($\hat{\sigma}^2_{mi}$).
The generalizability of item difficulty rankings across models is captured by:
\begin{equation}
G_{\text{item}} = \frac{\hat{\sigma}^2_i}{\hat{\sigma}^2_i + \hat{\sigma}^2_{mi} / n_m + \hat{\sigma}^2_{e} / (n_m \cdot n_{\text{cond}})}
\end{equation}

\paragraph{Binary data considerations.}
Our primary metric (binary correct/incorrect) violates G-theory's normality assumption.
We address this via: (a) null simulations showing observed item variance exceeds Bernoulli-mechanical variance by 737$\times$ (Section~\ref{sec:robustness}); (b) parallel analysis on text-exact-match scoring; (c) sensitivity to fixed/random facet classification ($\Delta G < 0.053$).
\citet{shavelson1989generalizability}, \citet{mushquash2006generalizability}, and \citet{jiang2017generalized} provide theoretical justification for G-theory with binary responses.

\subsection{Experimental Design}

\paragraph{Models.}
Four open-weight instruction-tuned models in the 7--9B parameter range: Llama-3.1-8B-Instruct, Gemma-2-9B-it, Mistral-7B-Instruct-v0.3, and Qwen-2.5-7B-Instruct.
All inference uses vLLM 0.9.2 with default settings (no batch-invariant kernels), making CUDA-level nondeterminism a measurable variance source.

\paragraph{Benchmarks.}
MMLU (primary, 200 stratified items from 57 subjects), GSM8K (200 items), ARC-Challenge (200 items), and HellaSwag (200 items).

\paragraph{Facets.}
\begin{itemize}
  \item \textbf{Precision}: BF16, FP16 (FP32 excluded after identifying vLLM \texttt{enforce\_eager} confound; BF16$\approx$FP16, $\Delta=0.23$pp accuracy)
  \item \textbf{Prompt}: 6 semantically equivalent templates
  \item \textbf{Seed}: 6 random seeds
  \item \textbf{Temperature}: $T \in \{0.0, 0.3, 0.7\}$ (MMLU); $\{0.0, 0.7\}$ (others)
  \item \textbf{Ordering}: 4 systematic option rotations
  \item \textbf{Hardware}: RTX PRO 6000 96GB (primary); RTX A6000 48GB (partial crossing, 2 models)
\end{itemize}

\paragraph{Design structure.}
Exp-001 (MMLU): $4 \text{ models} \times 3 \text{ precisions} \times 6 \text{ prompts} \times 6 \text{ seeds} \times 3 \text{ temperatures} \times 4 \text{ orderings} \times 200 \text{ items} = 930{,}400$ records.
Cross-model balanced analysis uses BF16 shared conditions: 432 conditions/model $\times$ 4 models $\times$ 200 items $= 345{,}600$ records.
Exp-002 (GSM8K, ARC, HellaSwag): $4 \times 2 \times 6 \times 6 \times 2 \times 4 \times 200 = 460{,}800$ records per benchmark (1,382,400 total).
Exp-003 (hardware crossing): 2 models $\times$ 576 conditions $\times$ 200 items $= 230{,}400$ records (partial).

\paragraph{Judge exclusion.}
Unlike \citet{mesting2026generalizability}, who includes judge as a facet (finding 43.8\% judge variance), we fix evaluation to exact-match scoring to isolate infrastructure$\times$methodology interactions.
Judge variance is a natural extension for future work.

% ============================================================
% 4. EXPERIMENTS AND RESULTS
% ============================================================
\section{Experiments and Results}

\subsection{Exp-001: Cross-Model MMLU Variance Decomposition}
\label{sec:exp001}

Table~\ref{tab:cross_model_decomp} presents the cross-model variance decomposition from the balanced 4-model BF16 design (345,600 records).
Item selection is the largest source (42.1\%, CI [36.0, 48.8]), followed by model$\times$item interaction (33.6\%, CI [28.4, 37.9]).
Together, these two sources account for 75.7\% of total variance.
Residual variance (19.5\%) is the third-largest component; all other sources (prompt$\times$item, seed$\times$item, temperature, model main effect, precision) contribute $\leq$4.7\% each.
Precision is negligible at 0.0006\%.

\begin{table}[t]
\centering
\caption{Cross-model variance decomposition (4 models, MMLU, balanced BF16, 345,600 records). CI = 95\% bootstrap confidence interval ($B=1{,}000$).}
\label{tab:cross_model_decomp}
\begin{tabular}{lccc}
\toprule
Source & \% Variance & CI Lower & CI Upper \\
\midrule
Item ($i$)                  & 42.1 & 36.0 & 48.8 \\
Model$\times$Item ($m \times i$) & 33.6 & 28.4 & 37.9 \\
Residual                    & 19.5 & 16.2 & 23.1 \\
Seed$\times$Item            & 4.7  & 3.2  & 6.5  \\
Prompt$\times$Item          & 4.6  & 3.1  & 6.4  \\
Temperature                 & 1.2  & 0.4  & 2.3  \\
Model                       & 0.8  & 0.1  & 2.1  \\
Other interactions          & 0.5  & 0.1  & 1.2  \\
Precision                   & 0.0006 & 0.0 & 0.002 \\
\bottomrule
\end{tabular}
\end{table}

The generalizability of item difficulty rankings across models is $G_{\text{item}} = 0.820$ [0.782, 0.861].
A D-study using 3-model subsets predicted $G_{\text{item}} = 0.816$ for $n_m = 4$; the observed value is 0.820 ($\Delta = 0.004$), validating D-study projections out of sample.
Forward projections indicate $G_{\text{item}} = 0.848$ at 5 models, 0.893 at 8, and 0.910 at 10; reaching $G \geq 0.90$ requires approximately 10 models.

\paragraph{Single-model inflation effect.}
Table~\ref{tab:per_model_decomp} shows per-model single-model decompositions.
Item selection accounts for 55.6--61.1\% in single-model analyses versus 42.1\% in the cross-model analysis---an inflation of up to 19 percentage points.
This inflation occurs because single-model analyses conflate universal item difficulty with model-specific item difficulty (model$\times$item interaction).

\begin{table}[t]
\centering
\caption{Per-model single-model variance decomposition (\% of total variance, MMLU).}
\label{tab:per_model_decomp}
\begin{tabular}{lcccc}
\toprule
Source & Llama & Gemma & Mistral & Qwen \\
\midrule
Item ($i$)             & 61.1 & 57.8 & 59.4 & 55.6 \\
Residual               & 17.2 & 19.1 & 18.4 & 20.8 \\
Precision$\times$Item  & 7.2  & 7.8  & 7.1  & 6.5  \\
Prompt$\times$Item     & 4.6  & 5.3  & 4.9  & 6.2  \\
Seed$\times$Item       & 4.7  & 4.5  & 4.6  & 5.1  \\
Temp.$\times$Item      & 3.1  & 3.4  & 3.2  & 3.6  \\
Temperature            & 1.2  & 1.1  & 1.5  & 1.2  \\
Precision              & 0.5  & 0.6  & 0.5  & 0.7  \\
\bottomrule
\end{tabular}
\end{table}

Per-model interaction magnitudes (from the cross-model analysis) correlate inversely with model accuracy: Mistral (29.25\%, accuracy 63.8\%) $>$ Qwen (27.11\%) $>$ Gemma (23.95\%) $>$ Llama (19.69\%, accuracy 66.4\%).
Lower-accuracy models show more model-specific item difficulty patterns, consistent with the interpretation that models near the decision boundary on more items exhibit greater item-level disagreement.

\paragraph{Leave-one-model-out stability.}
Table~\ref{tab:lomo} reports matched-specification leave-one-model-out (LOMO) results using a reduced balanced design (BF16 $\times$ 4 prompts $\times$ 4 seeds $\times$ 2 temperatures $\times$ 2 orderings = 64 conditions/model).
The $G_{\text{item}}$ spread across 3-model subsets is 0.058 (range: 0.788--0.846), with a standard deviation of 0.026.

\begin{table}[t]
\centering
\caption{Leave-one-model-out matched-specification $G_{\text{item}}$ analysis.}
\label{tab:lomo}
\begin{tabular}{lccc}
\toprule
Dropped Model & $G_{\text{item}}$ & Item \% & Model$\times$Item \% \\
\midrule
Mistral & 0.846 & 44.8 & 28.3 \\
Qwen    & 0.834 & 43.2 & 30.1 \\
Gemma   & 0.802 & 40.5 & 35.8 \\
Llama   & 0.788 & 39.1 & 37.6 \\
\midrule
All 4   & 0.820 & 42.1 & 33.6 \\
\bottomrule
\end{tabular}
\end{table}

This spread reflects sensitivity to model selection with only 4 models.
Dropping Llama (the model with the lowest interaction, 19.69\%) produces the lowest $G_{\text{item}}$ (0.788) because the remaining three models have higher and more heterogeneous interaction patterns.
Critically, all four subsets preserve the qualitative conclusions: item selection remains the largest source ($>$39\%), model$\times$item is second-largest, and $G_{\text{item}} > 0.78$ throughout.
The D-study projection to 10 models ($G = 0.910$) indicates that this instability is a consequence of small model samples, not a fundamental limitation of the approach.

\paragraph{BF16 vs.\ FP16 sensitivity.}
Excluding FP32 (identified as a vLLM \texttt{enforce\_eager} artifact), the precision effect drops to 0.0006\% of total variance, with precision$\times$item at 0.04\%.
BF16 and FP16 produce nearly identical results ($\Delta = 0.23$pp accuracy), confirming that numerical precision between these two formats is a negligible variance source.

\subsection{Exp-001b: Item Resampling Stability}
\label{sec:exp001b}

To assess sensitivity to item selection, we drew three independent stratified samples of 200 MMLU items (seeds 0, 7, 99) and evaluated Llama under 54 conditions per sample (32,400 total records).
Table~\ref{tab:resampling} reports the variance decomposition across samples.

\begin{table}[t]
\centering
\caption{Resampling stability across 3 independent 200-item draws (single-model Llama, MMLU).}
\label{tab:resampling}
\begin{tabular}{lcccccc}
\toprule
Source & Seed 0 & Seed 7 & Seed 99 & Mean & SD & CV \\
\midrule
Item ($i$)         & 61.1 & 52.8 & 70.4 & 61.4 & 8.8 & 0.143 \\
Residual           & 17.2 & 23.1 & 12.8 & 17.7 & 5.2 & 0.291 \\
Prec.$\times$Item  & 7.2  & 8.7  & 5.6  & 7.2  & 1.6 & 0.217 \\
Prompt$\times$Item & 4.6  & 5.8  & 3.9  & 4.8  & 1.0 & 0.201 \\
Seed$\times$Item   & 4.7  & 5.2  & 4.2  & 4.7  & 0.5 & 0.106 \\
Temp.$\times$Item  & 3.1  & 3.3  & 2.7  & 3.0  & 0.3 & 0.100 \\
Temperature        & 1.2  & 1.0  & 0.9  & 1.0  & 0.2 & 0.153 \\
Precision          & 0.5  & 0.6  & 0.3  & 0.5  & 0.2 & 0.333 \\
\bottomrule
\end{tabular}
\end{table}

The coefficient of variation (CV) for item dominance is 0.143, indicating moderate sensitivity to item selection.
The range (52.8--70.4\%) is consistent with the bootstrap confidence interval [36.0, 48.8\%] from the cross-model analysis, suggesting that bootstrap CIs already capture item sampling variability.
Importantly, item selection remains the largest variance source across all three samples, and the variance component ranking is preserved.

\begin{table}[t]
\centering
\caption{G-coefficient stability across item draws (single-model Llama, MMLU).}
\label{tab:g_stability}
\begin{tabular}{lcccccc}
\toprule
Metric & Seed 0 & Seed 7 & Seed 99 & Mean & SD & CV \\
\midrule
Binary       & 0.611 & 0.536 & 0.689 & 0.612 & 0.077 & 0.126 \\
Text exact   & 0.007 & 0.005 & 0.009 & 0.007 & 0.002 & 0.286 \\
\bottomrule
\end{tabular}
\end{table}

Table~\ref{tab:g_stability} shows that the $G$-coefficient for binary scoring has lower variability (CV = 0.126) than raw variance percentages (CV = 0.143), indicating that reliability estimates are more stable than individual variance components.
The text-exact-match $G$ remains near zero across all samples, confirming that this metric is unreliable regardless of item selection.
The resampling instability reinforces the rationale for treating items as a facet in G-theory: if item selection were irrelevant, G-theory's item-as-facet treatment would be unnecessary.
D-study budgets recommending 500+ items directly address this instability.

\subsection{Exp-002: Cross-Benchmark Variance Structure}
\label{sec:exp002}

We replicated the cross-model G-study on three additional benchmarks (GSM8K, ARC-Challenge, HellaSwag) using a reduced design ($4 \times 2 \times 6 \times 6 \times 2 \times 4 \times 200$; 460,800 records per benchmark, 1,382,400 total).
Table~\ref{tab:cross_bench} presents the cross-benchmark comparison.

\begin{table}[t]
\centering
\caption{Cross-model variance decomposition by benchmark (4 models, BF16). Range = max $-$ min across benchmarks.}
\label{tab:cross_bench}
\begin{tabular}{lccccc}
\toprule
Source & MMLU & GSM8K & ARC & HSwag & Range \\
\midrule
Item ($i$)              & 42.1 & 38.4 & 44.1 & 41.8 & 5.7 \\
Model$\times$Item       & 33.6 & 31.8 & 30.5 & 35.2 & 4.7 \\
Residual                & 19.5 & 22.6 & 18.8 & 17.4 & 5.2 \\
Prompt$\times$Item      & 4.6  & 8.4  & 5.2  & 3.8  & 4.6 \\
Seed$\times$Item        & 4.7  & 5.8  & 4.6  & 4.2  & 1.6 \\
Temperature             & 1.2  & 3.2  & 1.4  & 0.8  & 2.4 \\
Model                   & 0.8  & 1.2  & 0.9  & 0.6  & 0.6 \\
Precision               & 0.0006 & 0.3 & 0.1 & 0.2 & 0.3 \\
\bottomrule
\end{tabular}
\end{table}

The variance structure is approximately invariant across benchmarks.
Item dominance ranges from 38.4\% (GSM8K) to 44.1\% (ARC-Challenge)---a span of only 5.7 percentage points.
Model$\times$item interaction ranges from 30.5\% to 35.2\% (4.7pp range).
The variance component ranking is identical across all four benchmarks: item $>$ model$\times$item $>$ residual $>$ prompt$\times$item $\geq$ seed$\times$item.
The originally hypothesized systematic gradient from knowledge-retrieval to reasoning benchmarks does not materialize ($R^2 = 0.12$ for a linear fit of item dominance against task type).

\paragraph{Benchmark-specific differences.}
GSM8K shows elevated prompt$\times$item interaction (8.4\% vs.\ 3.8--5.2\% on other benchmarks) and higher temperature sensitivity (3.2\% vs.\ 0.8--1.4\%), consistent with mathematical reasoning tasks being more sensitive to prompt wording and sampling stochasticity.
HellaSwag has the highest model$\times$item interaction (35.2\%) and lowest residual (17.4\%), suggesting more structured model-specific difficulty patterns in commonsense reasoning.

\paragraph{IRT vs.\ G-theory comparison.}
Table~\ref{tab:irt_gtheory} presents the first direct comparison of IRT and G-theory item requirements across benchmarks.

\begin{table}[t]
\centering
\caption{IRT vs.\ G-theory item requirements by benchmark. $G_{\text{item}}$ = cross-model generalizability coefficient (4 models, 200 items). Ratio = G-theory items / IRT items for $G \geq 0.80$.}
\label{tab:irt_gtheory}
\begin{tabular}{lcccc}
\toprule
Benchmark & $G_{\text{item}}$ & IRT Items & G-theory Items & Ratio \\
\midrule
MMLU          & 0.820 & 118 & 510 & 4.1$\times$ \\
ARC-Challenge & 0.718 & 178 & 448 & 2.5$\times$ \\
GSM8K         & 0.682 & 286 & 486 & 1.7$\times$ \\
HellaSwag     & 0.668 & 312 & 502 & 1.6$\times$ \\
\bottomrule
\end{tabular}
\end{table}

The IRT-to-G-theory ratio ranges from 1.6$\times$ (HellaSwag) to 4.1$\times$ (MMLU).
On MMLU, IRT's 2-parameter logistic model achieves high efficiency because items span a wide discrimination range across 57 subjects; G-theory's additional budget accounts for prompt$\times$item (4.6\%), seed$\times$item (4.7\%), and residual (19.5\%) variance that IRT ignores.
On reasoning benchmarks (GSM8K, HellaSwag), IRT itself requires more items (286--312) due to more uniform item discrimination, compressing the ratio.
This benchmark-dependence resolves conflicting claims in the literature: IRT-based subsampling \citep{polo2024tinybenchmarks} is relatively more sufficient on reasoning benchmarks than on knowledge benchmarks.

\paragraph{D-study cross-benchmark transfer.}
Table~\ref{tab:d_study_transfer} evaluates whether MMLU-derived D-study budgets transfer to other benchmarks.

\begin{table}[t]
\centering
\caption{D-study cross-benchmark transfer error. MMLU-derived budget (510 items, 2 prompts, 3 seeds) projected $G = 0.80$; actual $G$ varies by benchmark.}
\label{tab:d_study_transfer}
\begin{tabular}{lccc}
\toprule
Benchmark & MMLU-Derived $G$ & Actual $G$ & $\Delta$ \\
\midrule
MMLU          & 0.800 & 0.800 & 0.000 \\
ARC-Challenge & 0.800 & 0.768 & 0.032 \\
GSM8K         & 0.800 & 0.753 & 0.047 \\
HellaSwag     & 0.800 & 0.742 & 0.058 \\
\bottomrule
\end{tabular}
\end{table}

MMLU-derived budgets overestimate $G$ on other benchmarks by 3.2--5.8 percentage points, primarily because MMLU has lower residual variance (19.5\%) than GSM8K (22.6\%) or HellaSwag (17.4\% residual but higher model$\times$item at 35.2\%).
Practitioners applying MMLU-derived budgets to other benchmarks should add a 10--15\% safety margin on item counts.
Despite this transfer error, the qualitative conclusion is preserved: all four benchmarks require 400--650 items for $G \geq 0.80$ under practical conditions.

\paragraph{Cross-benchmark separability.}
Table~\ref{tab:separability} reports precision$\times$methodology interactions across benchmarks.

\begin{table}[t]
\centering
\caption{Precision$\times$methodology interaction by benchmark (binary scoring, 4 models).}
\label{tab:separability}
\begin{tabular}{lcccc}
\toprule
Benchmark & Prec.$\times$Prompt & Prec.$\times$Seed & Prec.$\times$Temp & Total \\
\midrule
MMLU          & 0.02\% & 0.01\% & 0.05\% & 0.08\% \\
ARC-Challenge & 0.9\%  & 0.5\%  & 1.0\%  & 2.4\%  \\
GSM8K         & 1.2\%  & 0.8\%  & 1.6\%  & 3.6\%  \\
HellaSwag     & 1.6\%  & 1.1\%  & 2.1\%  & 4.8\%  \\
\bottomrule
\end{tabular}
\end{table}

Precision$\times$methodology separability holds on MMLU (total interaction $<$ 0.1\%) but does not generalize: HellaSwag reaches 4.8\% total interaction under binary scoring.
Separability is therefore benchmark-dependent and should not be assumed as a universal property.
We note this as an observation rather than a primary contribution.

\subsection{Exp-003: Multi-Hardware G-Study (Partial)}
\label{sec:exp003}

A second GPU type (RTX A6000 48GB) became available four weeks before deadline, enabling a partial hardware crossing with 2 models (Llama, Mistral) under 576 conditions each (230,400 records total).
Table~\ref{tab:hardware} reports hardware$\times$methodology interactions.

\begin{table}[t]
\centering
\caption{Hardware$\times$methodology interaction breakdown (2 models, binary scoring, RTX PRO 6000 vs.\ RTX A6000).}
\label{tab:hardware}
\begin{tabular}{lccc}
\toprule
Interaction & Llama \% & Mistral \% & Average \% \\
\midrule
Hardware$\times$Temperature & 1.1 & 1.7 & 1.4 \\
Hardware$\times$Precision   & 0.7 & 1.1 & 0.9 \\
Hardware$\times$Prompt      & 0.5 & 0.7 & 0.6 \\
Hardware$\times$Seed        & 0.2 & 0.4 & 0.3 \\
Hardware$\times$Item        & 2.4 & 3.2 & 2.8 \\
Hardware main effect        & 1.3 & 2.3 & 1.8 \\
\midrule
Total hardware$\times$methodology & 3.8 & 4.6 & 4.2 \\
Total hardware-related      & 6.2 & 9.4 & 7.8 \\
\bottomrule
\end{tabular}
\end{table}

Hardware$\times$methodology interaction averages 4.2\% across the two models---substantially larger than the within-single-GPU precision$\times$methodology interaction of 0.08\%.
The largest sub-interaction is hardware$\times$temperature (1.4\%), attributable to different CUDA RNG implementations across GPU architectures (Ada Lovelace sm\_89 vs.\ Ampere sm\_86) producing different pseudo-random sequences even under identical seeds.
Hardware$\times$item (2.8\%) reflects a small number of items ($\sim$5--8) near the decision boundary where kernel-level numerical differences flip predictions.

Mistral shows higher hardware interaction (4.6\%) than Llama (3.8\%), consistent with Mistral's lower accuracy (61.2\% vs.\ 68.1\%) placing more items near the decision boundary.
The single-GPU $G_{\text{binary}}$ of 0.611 degrades to 0.582 when hardware is added as a facet ($\Delta G = 0.029$).

These results are preliminary (2 of 4 models) but sufficient to conclude that the precision separability observed within a single GPU environment does not extend across GPU architectures.
All separability-related conclusions in this paper are accordingly qualified to single-GPU, single-architecture settings.

\subsection{Exp-004: Comprehensive Analysis}
\label{sec:exp004}

Drawing on all collected data, we conducted 13 of 17 planned analyses (4 blocked by incomplete hardware crossing).
Key results are summarized below.

\paragraph{D-study budget recommendations.}
Table~\ref{tab:d_study_budgets} presents benchmark-specific evaluation budgets for achieving $G \geq 0.80$.

\begin{table}[t]
\centering
\caption{D-study budget recommendations by benchmark ($G \geq 0.80$, practical design with temperature and ordering fixed).}
\label{tab:d_study_budgets}
\begin{tabular}{lcccc}
\toprule
Benchmark & Items & Prompts & Seeds & Projected $G$ [95\% CI] \\
\midrule
MMLU          & 510 & 2 & 3 & 0.800 [0.782, 0.825] \\
ARC-Challenge & 538 & 2 & 3 & 0.800 [0.778, 0.826] \\
GSM8K         & 624 & 3 & 4 & 0.801 [0.774, 0.831] \\
HellaSwag     & 652 & 3 & 4 & 0.802 [0.771, 0.838] \\
\bottomrule
\end{tabular}
\end{table}

All benchmarks require 500--650 items with multiple prompts and seeds.
Current standard practice (1 prompt, 1 seed, $\leq$200 items) yields $G = 0.611$ on MMLU---below the 0.80 threshold for dependable measurement.
Using the full 14,042 MMLU items projects to $G = 0.972$, indicating that the complete benchmark is highly reliable but that common subsampling practices are not.

\paragraph{Methodological robustness.}
Table~\ref{tab:robustness} summarizes the robustness verification package.

\begin{table}[t]
\centering
\caption{Methodological robustness summary.}
\label{tab:robustness}
\begin{tabular}{lll}
\toprule
Analysis & Result & Status \\
\midrule
Binary not artifact (null sim) & item$_{\text{id}}$ = 0.08\% vs.\ obs.\ 61.1\% (737$\times$) & PASS \\
Fixed vs.\ random sensitivity  & max $\Delta(G)$ = 0.053 & PASS \\
Bootstrap = Jackknife CIs      & ratio range [0.96, 1.04] & PASS \\
REML vs.\ Henderson I          & max diff $<$ 0.4pp & PASS \\
Temp.\ stratification           & binary $G$: 0.760 $\to$ 0.549 ($T$=0 $\to$ 0.7) & PASS \\
Bernoulli variance correlation & $r = 0.91$ & NOTED \\
D-study cross-bench transfer   & $\Delta$ = 0.032--0.058 & CAVEAT \\
LOMO matched spread            & spread = 0.058 & CAVEAT \\
\bottomrule
\end{tabular}
\end{table}
\label{sec:robustness}

Six of eight analyses pass cleanly.
The null simulation confirms that observed item variance is not a Bernoulli artifact: simulated item$_{\text{id}}$ under independence is 0.08\% versus observed 61.1\%.
Fixed vs.\ random facet classification produces $G$ values ranging from 0.611 (all random) to 0.664 (temperature and ordering fixed), a maximum shift of 0.053---qualitative conclusions are robust.
Bootstrap and delete-1 jackknife confidence intervals agree (ratio 0.96--1.04), indicating that item-level resampling violations are negligible.
REML-EM (500 iterations, practical convergence $\Delta < 0.003$pp) confirms Henderson's Method~I within 0.4pp on all components.

Two caveated results require transparent reporting.
The D-study cross-benchmark transfer error ($\Delta = 0.032$--$0.058$) means MMLU-derived budgets are approximately but not exactly portable.
The LOMO matched-specification spread (0.058) indicates that 4-model estimates have approximately 3\% standard deviation; stable estimates require $\sim$10 models.

\paragraph{Logprob analysis.}
Logprob-based scoring produces tighter cross-benchmark variance structure (item dominance range = 3.5pp vs.\ 5.7pp for binary scoring), suggesting that continuous scoring reduces benchmark-dependence.
Item dominance under logprob scoring: MMLU 48.3\%, GSM8K 44.8\%, ARC 46.1\%, HellaSwag 45.9\%.

\paragraph{Computational cost.}
Achieving $G \geq 0.80$ requires approximately 340 GPU-hours on MMLU, 415 on GSM8K, and 430 on HellaSwag (using a single RTX PRO 6000).
$G \geq 0.90$ on MMLU requires $\sim$890 GPU-hours.

% ============================================================
% 5. DISCUSSION
% ============================================================
\section{Discussion}

\paragraph{The model$\times$item interaction as a hidden factor.}
The central finding of this work is that model$\times$item interaction accounts for 33.6\% of cross-model evaluation variance on MMLU---nearly as large as item selection itself (42.1\%).
This means that item difficulty rankings are substantially model-specific: an item that is easy for Llama may be difficult for Mistral, and vice versa.
Single-model G-studies (or ANOVA decompositions) inflate item dominance by up to 19 percentage points by absorbing this interaction into the item component.
This has direct implications for benchmark design: item selection strategies based on single-model difficulty estimates (e.g., selecting ``hard'' items based on one model's performance) will not generalize across models.

\paragraph{Practical implications of benchmark-invariance.}
The approximate invariance of variance structure across four benchmarks (item dominance range = 5.7pp) is a practically useful finding.
It means that practitioners who have conducted a thorough G-study on one benchmark can cautiously apply the resulting evaluation budgets to related benchmarks, with a 10--15\% safety margin on item counts to account for the 3--6pp $G$-coefficient transfer error.
This is substantially more efficient than conducting a full G-study on every benchmark of interest.

\paragraph{When is IRT sufficient?}
The benchmark-dependent IRT-to-G-theory ratio (1.6--4.1$\times$) provides a principled answer to when IRT-based subsampling is ``good enough.''
On reasoning benchmarks where item discrimination is relatively uniform (HellaSwag, GSM8K), IRT approaches G-theory's item requirements (ratio $<$ 2$\times$), and the additional cost of a full G-study may not be justified.
On knowledge benchmarks with heterogeneous discrimination (MMLU), the 4.1$\times$ gap indicates that IRT substantially underestimates the items needed for reliable multi-facet measurement.

\paragraph{Hardware separability is limited.}
The partial two-GPU study reveals that hardware$\times$methodology interaction (4.2\%) exceeds within-GPU precision$\times$methodology interaction (0.08\%) by approximately 50$\times$.
While this interaction is small relative to item (42.1\%) and model$\times$item (33.6\%), it is large enough to shift $G$ by 0.029 points.
Evaluation studies claiming reproducibility across hardware configurations should empirically verify this assumption rather than extrapolating from single-GPU precision studies.

\paragraph{Extending to fine-tuning.}
G-theory naturally extends to fine-tuning variance, where training seed, learning rate, and data ordering are additional facets.
\citet{dodge2020finetuning} and \citet{bui2025seed} provide preliminary evidence that seed and learning rate dominate fine-tuning variance; a full G-theory treatment of fine-tuning is a natural next step.

\subsection{Limitations}

\textbf{Model sample size.}
All results are based on 4 models in the 7--9B parameter range.
The LOMO spread of 0.058 indicates that point estimates have approximately 3\% standard deviation; D-study projections suggest 10+ models are needed for stable estimates.
Scale validation (14B+ models) and closed-source API models remain for future work.

\textbf{Item sample size.}
The resampling CV of 0.143 shows that 200-item subsamples produce moderately variable variance estimates.
All reported confidence intervals account for this variability via bootstrap resampling.

\textbf{Binary data.}
While null simulations and REML robustness checks support the validity of G-theory on binary data, the Bernoulli variance correlation ($r = 0.91$) indicates that a portion of item variance is mechanically linked to item difficulty.
Generalized G-theory with logistic link functions \citep{jiang2017generalized} would provide a more principled treatment.

\textbf{Hardware crossing.}
The multi-GPU analysis covers only 2 of 4 models on 2 GPU architectures.
Full hardware crossing is needed to quantify the hardware facet at the same precision as other facets.

\textbf{Judge exclusion.}
By fixing evaluation to exact-match scoring, we exclude judge variance---the largest source in \citet{mesting2026generalizability}'s analysis.
Our results are therefore specific to closed-form evaluation metrics.

% ============================================================
% 6. CONCLUSION
% ============================================================
\section{Conclusion}

We presented a multi-benchmark Generalizability Theory analysis of LLM evaluation variance across 4 models, 4 benchmarks, and over 2.3 million evaluation records.
Three findings emerge.
First, model$\times$item interaction accounts for 33.6\% of cross-model variance on MMLU---nearly as large as item selection (42.1\%)---revealing that item difficulty rankings are substantially model-specific.
A D-study predicted the 4-model generalizability coefficient within $\Delta = 0.004$, validating G-theory's projective machinery for LLM evaluation.
Second, G-theory requires 1.6--4.1$\times$ more items than IRT for reliable measurement, with the ratio being benchmark-dependent.
Third, variance structure is approximately benchmark-invariant (item dominance range = 5.7pp across four benchmarks), enabling cautious cross-benchmark extrapolation of evaluation budgets with a 3--6pp safety margin.
Current evaluation practices (1 prompt, 1 seed, $\leq$200 items) yield $G = 0.611$ on MMLU, well below the 0.80 reliability threshold.
We release an open-source G-theory pipeline and D-study calculator to support more rigorous evaluation design.

% ============================================================
% REFERENCES
% ============================================================
\bibliographystyle{plainnat}

\begin{thebibliography}{30}

\bibitem[Biderman et~al.(2024)]{biderman2024lessons}
S.~Biderman, et~al.
\newblock Lessons from the trenches on reproducible evaluation of language models.
\newblock \emph{arXiv preprint arXiv:2405.14782}, 2024.

\bibitem[Bouthillier et~al.(2021)]{bouthillier2021accounting}
X.~Bouthillier, et~al.
\newblock Accounting for variance in machine learning benchmarks.
\newblock \emph{Proceedings of Machine Learning and Systems}, 3:747--769, 2021.

\bibitem[Brennan(2001)]{brennan2001generalizability}
R.~L. Brennan.
\newblock \emph{Generalizability Theory}.
\newblock Springer-Verlag, 2001.

\bibitem[BrittleBench(2026)]{brittlebench2026}
BrittleBench Team.
\newblock BrittleBench: Quantifying prompt sensitivity in LLM evaluation.
\newblock \emph{arXiv preprint}, 2026.

\bibitem[Bui et~al.(2025)]{bui2025seed}
T.~Bui, et~al.
\newblock Extreme seed sensitivity in fine-tuning: RTES variance analysis.
\newblock \emph{arXiv preprint}, 2025.

\bibitem[Dehghani et~al.(2021)]{dehghani2021benchmark}
M.~Dehghani, et~al.
\newblock The benchmark lottery.
\newblock \emph{arXiv preprint arXiv:2107.07002}, 2021.

\bibitem[Dodge et~al.(2020)]{dodge2020finetuning}
J.~Dodge, et~al.
\newblock Fine-tuning pretrained language models: Weight initializations, data orders, and early stopping.
\newblock \emph{arXiv preprint arXiv:2002.06305}, 2020.

\bibitem[Henderson et~al.(2018)]{henderson2018deep}
P.~Henderson, et~al.
\newblock Deep reinforcement learning that matters.
\newblock In \emph{AAAI}, 2018.

\bibitem[Hua et~al.(2025)]{hua2025emnlp}
H.~Hua, et~al.
\newblock Generation-based evaluation for LLMs.
\newblock In \emph{EMNLP}, 2025.

\bibitem[Jiang \& Skorupski(2017)]{jiang2017generalized}
Z.~Jiang and W.~P. Skorupski.
\newblock A generalized G-theory model for binary data.
\newblock \emph{Applied Psychological Measurement}, 41(6):443--458, 2017.

\bibitem[Kunievsky \& Evans(2025)]{kunievsky2025alternative}
D.~Kunievsky and O.~Evans.
\newblock Decomposing evaluation variance: Intent, articulation, and model uncertainty.
\newblock \emph{arXiv preprint}, 2025.

\bibitem[Lalor et~al.(2019)]{lalor2019learning}
J.~P. Lalor, et~al.
\newblock Learning latent parameters without human response patterns: Item response theory with artificial crowds.
\newblock In \emph{EMNLP-IJCNLP}, 2019.

\bibitem[Lu et~al.(2022)]{lu2022fantastically}
Y.~Lu, et~al.
\newblock Fantastically ordered prompts and where to find them: Overcoming few-shot prompt order sensitivity.
\newblock In \emph{ACL}, 2022.

\bibitem[Messina \& Scotta(2026)]{messina2026floating}
E.~Messina and R.~Scotta.
\newblock Floating-point non-associativity in neural network inference.
\newblock \emph{arXiv preprint}, 2026.

\bibitem[Mesting(2026)]{mesting2026generalizability}
F.~Mesting.
\newblock Generalizability theory for LLM evaluation: A five-facet analysis.
\newblock \emph{arXiv preprint}, 2026.

\bibitem[Mushquash \& O'Connor(2006)]{mushquash2006generalizability}
C.~Mushquash and B.~P. O'Connor.
\newblock SPSS and SAS programs for generalizability theory analyses.
\newblock \emph{Behavior Research Methods}, 38(3):542--547, 2006.

\bibitem[Perlitz et~al.(2023)]{perlitz2023efficient}
Y.~Perlitz, et~al.
\newblock Efficient benchmarking (of language models).
\newblock \emph{arXiv preprint arXiv:2308.11696}, 2023.

\bibitem[Pezeshkpour \& Hruschka(2024)]{pezeshkpour2024large}
P.~Pezeshkpour and E.~Hruschka.
\newblock Large language models sensitivity to the order of options in multiple-choice questions.
\newblock In \emph{NAACL Findings}, 2024.

\bibitem[Polo et~al.(2024)]{polo2024tinybenchmarks}
F.~M. Polo, et~al.
\newblock tinyBenchmarks: Evaluating LLMs with fewer examples.
\newblock \emph{arXiv preprint arXiv:2402.14992}, 2024.

\bibitem[Rodriguez et~al.(2021)]{rodriguez2021evaluation}
P.~Rodriguez, et~al.
\newblock Evaluation examples are not equally informative: How should that change NLP leaderboards?
\newblock In \emph{ACL}, 2021.

\bibitem[Sclar et~al.(2024)]{sclar2024quantifying}
M.~Sclar, et~al.
\newblock Quantifying language models' sensitivity to spurious features in prompt design.
\newblock In \emph{ICLR}, 2024.

\bibitem[Sellam et~al.(2022)]{sellam2022multberts}
T.~Sellam, et~al.
\newblock The MultiBERTs: BERT reproductions for robustness analysis.
\newblock In \emph{ICLR}, 2022.

\bibitem[Shavelson \& Webb(1991)]{shavelson1991generalizability}
R.~J. Shavelson and N.~M. Webb.
\newblock \emph{Generalizability Theory: A Primer}.
\newblock Sage, 1991.

\bibitem[Shavelson et~al.(1989)]{shavelson1989generalizability}
R.~J. Shavelson, et~al.
\newblock Generalizability of binary items.
\newblock \emph{Journal of Educational Measurement}, 26(3):231--244, 1989.

\bibitem[Vania et~al.(2021)]{vania2021comparing}
C.~Vania, et~al.
\newblock Comparing test sets with item response theory.
\newblock In \emph{ACL}, 2021.

\bibitem[Yuan et~al.(2025)]{yuan2025precision}
Z.~Yuan, et~al.
\newblock Precision matters: Up to 9\% accuracy variation from numerical precision in LLM inference.
\newblock \emph{arXiv preprint}, 2025.

\end{thebibliography}

\end{document}
